\section{Uniqueness of multi-block compressions}%
\label{sec:asym_uniqueness}

Theorem~\ref{thm:asym_multiblock} produces, for a tensor $B$ with the periodic vectors of a sum of
normal blocks, a block upper triangular form and compression maps $W_s$, $V_s$. This section
describes every such form of one tensor and the relation between any two of them. The description
is in terms of the labelled invariant flag of Definition~\ref{def:asym_flag_data}: the form
determines the flag, the flag determines the compression maps up to an explicit group, and two
flags of the same tensor are compared by an explicit matching. Throughout, a \emph{compression} of
$B$ onto the family $(C_s)_{s\in S}$ with $z$ zero slots is the datum of
Theorem~\ref{thm:asym_multiblock}, clauses (i)--(iii): an ordering $\sigma$ of the $r=|S|+z$ slots
and an invertible $G$ with $GB^iG^{-1}$ block upper triangular, diagonal block $C_s^i$ at the
position of $s$ and the $1\times1$ zero at a zero slot; the compression maps $W_b$ are the row
blocks of $G$ and $V_b$ the column blocks of $G^{-1}$, for every slot $b$, so that
$W_bV_{b'}=\delta_{bb'}1$. For a block slot $s$ the maps are those of
Theorem~\ref{thm:asym_multiblock}, and $W_sB^wV_s=C_s^w$ for every word $w$. The blocks are
assumed normal and of positive bond dimension, as in Theorem~\ref{thm:asym_multiblock}; the
arguments use of normality only that the module of each block is simple and has a nonzero
letter.

\subsection{Density and independence of traces}

A tensor $C=(C^i)_i$ of bond dimension $D$ is \emph{irreducible} when its letters have no common
invariant subspace other than $0$ and $\C^D$. An element $a=\sum_wc_ww$ of the word algebra $F$
acts on $C$ by $a(C)=\sum_wc_wC^w$, the empty word acting as $\Id$; $F_+$ is the ideal of the
combinations of nonempty words.

\begin{theorem}[The density theorem with targets]
    \label{thm:asym_uniq_density}
    \lean{MPSTensor.exists_eval_eq_of_pairwise_not_isGaugeRelated,
        MPSTensor.exists_aug_eq_zero_eval_eq_of_pairwise_not_isGaugeRelated}
    \leanok
    \uses{def:asym_word_algebra, def:asym_gauge_related}
    Let $C_1,\ldots,C_m$ be irreducible tensors of bond dimensions $D_1,\ldots,D_m$, pairwise not
    gauge related, and let $Y_j\in\MN{D_j}$ for every $j$. Then there is $a\in F$ with $a(C_j)=Y_j$
    for every $j$. If moreover every $C_j$ has a nonzero letter, $a$ can be chosen in $F_+$.
\end{theorem}

\begin{proof}
    \leanok
    \uses{thm:asym_separation}
    By Burnside's theorem the words act on $\C^{D_j}$ as the full matrix algebra, so there is
    $b_j\in F_+$ with $b_j(C_j)=Y_j$ whenever $C_j$ has a nonzero letter. The word-algebra modules
    of the $C_j$ are simple and pairwise non-isomorphic, and Theorem~\ref{thm:asym_separation} gives
    $p_j\in F_+$ with $p_j(C_k)=\delta_{jk}\Id$ for every $C_j$ with a nonzero letter. An
    irreducible tensor of positive bond dimension whose letters all vanish has $D=1$, and two such
    are gauge related, so there is at most one, $C_{j_0}$. Then
    $a=\sum_jb_jp_j+Y_{j_0}\bigl(1-\sum_jp_j\bigr)$, the second term present only when $C_{j_0}$
    exists, satisfies $a(C_j)=Y_j$ for every $j$, and lies in $F_+$ when there is no $C_{j_0}$.
\end{proof}

\begin{corollary}[Independence of the word traces]
    \label{cor:asym_uniq_trace_independence}
    \lean{MPSTensor.eq_zero_of_forall_sum_mul_trace_evalWord_eq_zero}
    \leanok
    \uses{thm:asym_uniq_density}
    Let $C_1,\ldots,C_m$ be irreducible tensors, pairwise not gauge related, each with a nonzero
    letter. If $c\in\C^m$ satisfies $\sum_jc_j\tr(C_j^w)=0$ for every nonempty word $w$, then $c=0$.
\end{corollary}

\begin{proof}
    \leanok
    \uses{thm:asym_uniq_density}
    The relation extends linearly to $\sum_jc_j\tr(a(C_j))=0$ for every $a\in F_+$.
    Theorem~\ref{thm:asym_uniq_density} gives $a_j\in F_+$ with $a_j(C_k)=\delta_{jk}\Id$, and the
    relation at $a_j$ reads $c_jD_j=0$. A tensor with a nonzero letter has $D_j>0$, so $c_j=0$.
\end{proof}

\subsection{Compressions and labelled flags}

\begin{proposition}[A compression is a labelled flag with an adapted basis]
    \label{prop:asym_uniq_dictionary}
    \lean{MPSTensor.MultiBlockCompression.flag,
        MPSTensor.MultiBlockCompression.toFlagData,
        MPSTensor.MultiBlockCompression.flag_covBy,
        MPSTensor.FlagData.toMultiBlockCompression,
        MPSTensor.FlagData.toMultiBlockCompression_flag}
    \leanok
    \uses{thm:asym_multiblock, def:asym_flag_data}
    Let $(\sigma,G)$ be a compression of $B$ and let $E_k\subseteq\C^{D_B}$ be the span of the
    first $k$ coordinate blocks. Then $M_k:=G^{-1}E_k$, $k=0,\ldots,r$, is a labelled invariant flag
    of the word-algebra module of $B$ over $(S,C)$ with $z$ zero steps, the step $M_k/M_{k-1}$
    carrying the slot $\sigma(k)$. Conversely every labelled invariant flag is the flag of a
    compression whose order is the order of the labels.
\end{proposition}

\begin{proof}
    \leanok
    \uses{def:asym_flag_data}
    Block upper triangular matrices map $E_k$ into $E_k$, so $B^iM_k=G^{-1}(GB^iG^{-1})E_k\subseteq M_k$;
    the letters induced on $E_k/E_{k-1}$ by $GB^iG^{-1}$ are the diagonal blocks, which are $C_s^i$
    on a block slot and zero on a one-dimensional zero slot. Conversely, choose a complement of
    $M_{k-1}$ in $M_k$ for every $k$ and a basis of it in which the induced letters are $C_s^i$
    (the label provides one); the matrix $G$ whose rows are the dual basis of the union of these
    bases is block upper triangular with the prescribed diagonal blocks. The flag of the
    compression so obtained is the given flag.
\end{proof}

\begin{lemma}[The compression maps in terms of the flag]
    \label{lem:asym_uniq_maps}
    \lean{MPSTensor.MultiBlockCompression.leftRow_mulVec_eq_zero_of_mem,
        MPSTensor.MultiBlockCompression.leftRow_surjOn,
        MPSTensor.MultiBlockCompression.leftRow_mulVec_mulVec_of_mem,
        MPSTensor.MultiBlockCompression.rightCol_mulVec_mem_flag,
        MPSTensor.MultiBlockCompression.leftRow_mul_rightCol_self}
    \leanok
    \uses{prop:asym_uniq_dictionary}
    Let $b$ be the slot at position $k$ of a compression with flag $(M_k)$ and let $C_b$ be its
    diagonal block ($C_s$ on a block slot, the $1\times1$ zero on a zero slot). Then $W_b$ vanishes
    on $M_{k-1}$, maps $M_k$ onto $\C^{n_k}$, and $W_bB^ix=C_b^iW_bx$ for $x\in M_k$; and $V_b$ maps
    $\C^{n_k}$ into $M_k$ with $W_bV_b=1$.
\end{lemma}

\begin{proof}
    \leanok
    For $x=G^{-1}e$ with $e\in E_k$, $W_bx$ is the $k$-th block of $e$, which vanishes exactly when
    $e\in E_{k-1}$; the $k$-th block of $(GB^iG^{-1})e$ is the diagonal block times the $k$-th block
    of $e$ because the row $k$ of a block upper triangular matrix meets only the columns $q\geq k$
    and $e$ has no component beyond $k$. The column block $V_b$ of $G^{-1}$ places a vector in the
    $k$-th block and $W_bV_b=1$ is $GG^{-1}=1$.
\end{proof}

\subsection{The matching of two flags}

\begin{lemma}[One step against another]
    \label{lem:asym_uniq_step}
    \lean{CoveringMatching.step_sup_eq_of_lt,
        CoveringMatching.step_sup_eq_of_lt',
        CoveringMatching.step_inf_eq_of_lt,
        CoveringMatching.step_lt_symm,
        MPSTensor.MultiBlockCompression.subquotEquivMatching}
    \leanok
    In a modular lattice let $P\lessdot Q$ and $P'\lessdot Q'$ be covering pairs and put
    $U=Q\wedge Q'$. If $P\vee(Q\wedge P')<P\vee(Q\wedge Q')$, then $P\vee U=Q$, $Q\wedge P'\leq P$,
    $P'\vee U=Q'$, $Q'\wedge P\leq P'$, $P\wedge Q'=Q\wedge P'$, and the strict inequality holds
    with the two pairs exchanged. For the lattice of invariant subspaces of a tensor, with
    $P\lessdot Q$ meaning that $Q/P$ is irreducible, the inclusions $U\subseteq Q$ and
    $U\subseteq Q'$ then induce isomorphisms $U/K\to Q/P$ and $U/K\to Q'/P'$ commuting with the
    letters, where $K=P\wedge Q'=Q\wedge P'$.
\end{lemma}

\begin{proof}
    \leanok
    $P\vee(Q\wedge P')$ and $P\vee U$ lie between $P$ and $Q$ with the first strictly below the
    second, so the first is $P$ and the second is $Q$, since nothing lies strictly between $P$ and
    $Q$. Then $P'\vee U$ equals $P'$ or $Q'$, and the first case gives $U\leq Q\wedge P'\leq P$, so
    $P\vee U=P\neq Q$; hence $P'\vee U=Q'$. Next $P'\vee(Q'\wedge P)$ equals $P'$ or $Q'$; in the
    second case $Q'\leq P\vee P'$ and the modular law gives
    $U\leq Q\wedge(P\vee P')=P\vee(Q\wedge P')=P$, again contradicting $P\vee U=Q$. The two
    inclusions $Q\wedge P'\leq P\wedge Q'$ and $Q'\wedge P\leq Q\wedge P'$ give the equality, and
    the exchanged inequality is $P'<Q'$. The isomorphisms are the second isomorphism theorem for
    $U$ against $P$ and against $P'$.
\end{proof}

\begin{proposition}[The matching]
    \label{prop:asym_uniq_matching}
    \lean{CoveringMatching.CoveringChain,
        CoveringMatching.CoveringChain.jump,
        CoveringMatching.CoveringChain.jump_spec,
        CoveringMatching.CoveringChain.jump_jump,
        CoveringMatching.CoveringChain.jumpEquiv,
        CoveringMatching.CoveringChain.length_eq,
        CoveringMatching.CoveringChain.jump_sup_eq,
        MPSTensor.MultiBlockCompression.matching,
        MPSTensor.MultiBlockCompression.matching_symm,
        MPSTensor.MultiBlockCompression.flag_sup_inf_eq,
        MPSTensor.MultiBlockCompression.flag_inf_eq,
        MPSTensor.MultiBlockCompression.isGaugeRelated_matching}
    \leanok
    \uses{lem:asym_uniq_step, prop:asym_uniq_dictionary}
    Let $(M_k)_{k=0}^r$ and $(\widetilde M_{k'})_{k'=0}^{\widetilde r}$ be the flags of two
    compressions of $B$. For every $k$ there is exactly one $k'$ with
    $M_{k-1}+(M_k\cap\widetilde M_{k'-1})\subsetneq M_{k-1}+(M_k\cap\widetilde M_{k'})$, namely the least
    $k'$ with $M_k\cap\widetilde M_{k'}\not\subseteq M_{k-1}$. Writing $k'=\pi(k)$, the map $\pi$ is a
    bijection, so $r=\widetilde r$, with inverse the same rule for the flags exchanged; and for
    every $k$ the letters induced on $M_k/M_{k-1}$ and on $\widetilde M_{\pi(k)}/\widetilde M_{\pi(k)-1}$
    are gauge related, both being isomorphic to the induced letters on $U/K$ with
    $U=M_k\cap\widetilde M_{\pi(k)}$ and $K=M_{k-1}\cap\widetilde M_{\pi(k)}=M_k\cap\widetilde M_{\pi(k)-1}$.
\end{proposition}

\begin{proof}
    \leanok
    \uses{lem:asym_uniq_step}
    The subspaces $M_{k-1}+(M_k\cap\widetilde M_j)$ increase with $j$ from $M_{k-1}$ to $M_k$, each
    equal to one of the two since $M_k/M_{k-1}$ is irreducible, so exactly one step is strict.
    Lemma~\ref{lem:asym_uniq_step} applied to the two covering pairs shows that the rule for the
    exchanged flags sends $k'$ back to $k$, and gives the identities and the isomorphisms.
\end{proof}

\begin{theorem}[The blocks are unique]
    \label{thm:asym_uniq_blocks}
    \lean{MPSTensor.MultiBlockCompression.z_eq,
        MPSTensor.MultiBlockCompression.labelMatching,
        MPSTensor.MultiBlockCompression.z_eq_and_exists_equiv_isGaugeRelated,
        MPSTensor.MultiBlockCompression.dim_labelMatching}
    \leanok
    \uses{prop:asym_uniq_matching, def:normal, def:asym_gauge_related}
    Let $B$ have compressions onto families $(C_s)_{s\in S}$ and $(\widetilde C_t)_{t\in\widetilde S}$
    of normal tensors of positive bond dimension, with $z$ and $\widetilde z$ zero slots. Then
    $z=\widetilde z$ and there is a bijection $\tau:S\to\widetilde S$ with $D_{\tau(s)}=D_s$ and
    $\widetilde C_{\tau(s)}$ gauge related to $C_s$.
\end{theorem}

\begin{proof}
    \leanok
    \uses{prop:asym_uniq_matching}
    The matching pairs steps with gauge related induced letters. On a zero step every letter is
    zero; on a block step some letter is nonzero, because a normal tensor of positive bond
    dimension has words of some positive length spanning a nonzero matrix algebra. Gauge related
    tensors have zero letters together, so zero steps are paired with zero steps, which gives
    $z=\widetilde z$, and block steps with block steps, which gives $\tau$.
\end{proof}

\begin{lemma}[The matching as a rank test]
    \label{lem:asym_uniq_rank_test}
    \lean{MPSTensor.MultiBlockCompression.leadingCols,
        MPSTensor.MultiBlockCompression.range_leadingCols,
        MPSTensor.MultiBlockCompression.range_fromCols_leadingCols,
        MPSTensor.MultiBlockCompression.rank_fromCols_leadingCols,
        MPSTensor.MultiBlockCompression.isLeast_rank_fromCols_leadingCols_eq,
        MPSTensor.MultiBlockCompression.coe_map_inf_eq_leadingCols}
    \leanok
    \uses{prop:asym_uniq_matching, prop:asym_uniq_dictionary}
    For two compressions $(\sigma,G)$ and $(\widetilde\sigma,\widetilde G)$ of $B$ with flags
    $(M_k)$ and $(\widetilde M_{k'})$, let $G_{\leq k}$ be the matrix of the columns of $G^{-1}$ at
    the slots of the positions $1,\ldots,k$, so that its column span is $M_k$, and define
    $\widetilde G_{\leq k'}$ likewise. Then the column span of the block matrix
    $(G_{\leq k}\ \widetilde G_{\leq k'})$ is $M_k+\widetilde M_{k'}$, so its rank is
    $\dim(M_k+\widetilde M_{k'})$; the subspace $M_k\cap\widetilde M_{k'}$ consists of the vectors
    $G_{\leq k}a$ with $G_{\leq k}a=\widetilde G_{\leq k'}a'$ for some $a'$; and, for normal blocks of
    positive bond dimension, the position $\pi(k)$ matched with $k$ is the least $k'$ with
    $\mathrm{rank}(G_{\leq k-1}\ \widetilde G_{\leq k'})=\mathrm{rank}(G_{\leq k}\ \widetilde G_{\leq k'})$.
\end{lemma}

\begin{proof}
    \leanok
    \uses{prop:asym_uniq_matching}
    The columns of $G^{-1}$ at the positions $1,\ldots,k$ span $G^{-1}E_k=M_k$. The two ranks agree
    exactly when $M_k\subseteq M_{k-1}+\widetilde M_{k'}$, that is, when
    $M_{k-1}+(M_k\cap\widetilde M_{k'})=M_k$ by the modular law; by
    Proposition~\ref{prop:asym_uniq_matching} the least such $k'$ is $\pi(k)$.
\end{proof}

\subsection{Uniqueness of the compression maps}

\begin{theorem}[Two compressions with the same flag]
    \label{thm:asym_uniq_same_flag}
    \lean{MPSTensor.MultiBlockCompression.withGauge,
        MPSTensor.MultiBlockCompression.blockTriangular_coordChange_of_flag_eq,
        MPSTensor.MultiBlockCompression.blockDiag'_coordChange_eq_smul,
        MPSTensor.MultiBlockCompression.coordChangeScalar_ne_zero,
        MPSTensor.MultiBlockCompression.leftRow_withGauge,
        MPSTensor.MultiBlockCompression.rightCol_withGauge,
        MPSTensor.MultiBlockCompression.conjTriangular,
        MPSTensor.MultiBlockCompression.conjTriangular_flag,
        MPSTensor.MultiBlockCompression.withGauge_eq_conjTriangular,
        MPSTensor.MultiBlockCompression.coordChange_conjTriangularGauge}
    \leanok
    \uses{lem:asym_uniq_maps, def:normal}
    Let $(\sigma,G)$ and $(\sigma,\widetilde G)$ be compressions of $B$ onto the same family with the
    same labelled flag, $\widetilde G^{-1}E_k=G^{-1}E_k$ for all $k$. Then $h=\widetilde GG^{-1}$ is
    block upper triangular for the decomposition of $\sigma$ and every diagonal block $h_{bb}$ is a
    nonzero scalar $\lambda_b$, so that
    \begin{align}
        \widetilde W_b & =\lambda_bW_b+\sum_{c\text{ after }b}h_{bc}W_c,                &
        \widetilde V_b & =\lambda_b^{-1}V_b+\sum_{c\text{ before }b}V_c\,(h^{-1})_{cb}.
    \end{align}
    Conversely, for every block upper triangular $h$ with nonzero scalar diagonal blocks,
    $(\sigma,hG)$ is a compression onto the same family with the same flag. The compressions with a
    given labelled flag are therefore a single orbit of the group of such matrices, acting freely.
\end{theorem}

\begin{proof}
    \leanok
    \uses{lem:asym_uniq_maps, lem:asym_normal_intertwiner}
    $h$ preserves every $E_k$, hence is block upper triangular, and so is $h^{-1}$. The diagonal
    block of $h(GB^iG^{-1})h^{-1}$ at a block slot $s$ is $h_{ss}C_s^ih_{ss}^{-1}$ and must equal
    $C_s^i$, so $h_{ss}$ commutes with every letter of the normal tensor $C_s$ and is a scalar; on a
    zero slot it is a nonzero number. The displayed formulas are the rows of $hG$ and the columns of
    $G^{-1}h^{-1}$. For the converse, conjugating a block upper triangular matrix by $h$ keeps it
    block upper triangular and keeps the diagonal blocks, and $(hG)^{-1}E_k=G^{-1}E_k$.
\end{proof}

\begin{theorem}[The compressions with one labelled flag form a torsor]
    \label{thm:asym_uniq_torsor}
    \lean{MPSTensor.MultiBlockCompression.bijOn_conjTriangular}
    \leanok
    \uses{thm:asym_uniq_same_flag}
    Let $(\sigma,G)$ be a compression of $B$ onto a family of normal tensors of positive bond
    dimension, and let $T_\sigma$ be the group of matrices that are block upper triangular for the
    decomposition of $\sigma$ with nonzero scalar diagonal blocks. Then $h\mapsto(\sigma,hG)$ is a
    bijection from $T_\sigma$ onto the set of compressions of $B$ onto the same family with the same
    number of zero slots, the same order and the same flag as $(\sigma,G)$.
\end{theorem}

\begin{proof}
    \leanok
    \uses{thm:asym_uniq_same_flag}
    By the converse part of Theorem~\ref{thm:asym_uniq_same_flag} the map takes values in the set;
    by its first part every compression $(\sigma,\widetilde G)$ of the set is $(\sigma,hG)$ with
    $h=\widetilde GG^{-1}\in T_\sigma$; and $h$ is recovered from $(\sigma,hG)$ as $(hG)G^{-1}$, so
    the map is injective.
\end{proof}

\begin{theorem}[Two compressions with different flags]
    \label{thm:asym_uniq_different_flags}
    \lean{MPSTensor.MultiBlockCompression.exists_intertwiner_of_matching,
        MPSTensor.MultiBlockCompression.existsUnique_intertwiner_of_matching,
        MPSTensor.MultiBlockCompression.exists_smul_of_matching_inl,
        MPSTensor.MultiBlockCompression.exists_smul_of_matching_inr,
        MPSTensor.MultiBlockCompression.exists_ne_zero_eq_smul_of_slotTensor_eq,
        MPSTensor.MultiBlockCompression.inf_eq_and_sup_eq_of_flag_eq}
    \leanok
    \uses{prop:asym_uniq_matching, lem:asym_uniq_maps}
    Let two compressions of $B$ onto families $(C_s)_{s\in S}$ and
    $(\widetilde C_t)_{t\in\widetilde S}$ of normal tensors of positive bond dimension have flags
    $(M_k)$ and $(\widetilde M_{k'})$ and matching $\pi$. Let $k$ carry the slot $b$ of the first
    compression and $k'=\pi(k)$ the slot $b'$ of the second, with diagonal blocks $C_b$ and
    $\widetilde C_{b'}$ (the $1\times1$ zero on a zero slot), and put $U=M_k\cap\widetilde M_{k'}$
    and $\mathcal N=M_{k-1}+\widetilde M_{k'-1}$. Then there is a unique invertible intertwiner $X$
    from $C_b$ to $\widetilde C_{b'}$ with
    \begin{align}
        \widetilde W_{b'}x             & =XW_bx\quad\text{for all }x\in U,    &
        (\widetilde V_{b'}-V_bX^{-1})v & \in\mathcal N\quad\text{for all }v .
    \end{align}
    On two zero slots $X$ is a nonzero scalar, and so it is when $\widetilde C_{b'}=C_b$, in
    particular for two compressions onto one family at two slots carrying the same label. When the
    two labelled flags coincide, $U=M_k$ and $\mathcal N=M_{k-1}$.
\end{theorem}

\begin{proof}
    \leanok
    \uses{prop:asym_uniq_matching, lem:asym_uniq_maps}
    By Proposition~\ref{prop:asym_uniq_matching}, $M_{k-1}+U=M_k$, $\widetilde M_{k'-1}+U=\widetilde M_{k'}$
    and $U\cap M_{k-1}=K=U\cap\widetilde M_{k'-1}$. By Lemma~\ref{lem:asym_uniq_maps}, $W_b$
    restricted to $U$ intertwines the letters with $C_b$, is onto $\C^{n_k}$ since
    $W_b(U)=W_b(M_{k-1}+U)=W_b(M_k)$, and has kernel $K$; the same for $\widetilde W_{b'}$. Define
    $Xv=\widetilde W_{b'}u$ for any $u\in U$ with $W_bu=v$; the choice does not matter since
    $\widetilde W_{b'}$ kills $K$. Then $X$ is bijective and intertwines, and it is unique because
    $W_b$ maps $U$ onto $\C^{n_k}$. For the injections, decompose $\widetilde V_{b'}v=y+u$ with
    $y\in\widetilde M_{k'-1}$, $u\in U$ and $V_bX^{-1}v=y'+u'$ with $y'\in M_{k-1}$, $u'\in U$; applying
    $\widetilde W_{b'}$ and $W_b$ gives $W_b(u-u')=0$, so $u-u'\in K\subseteq M_{k-1}$, and the
    difference lies in $\mathcal N$. On a zero slot $X$ is $1\times1$; when the two diagonal blocks
    are the same normal tensor, $X$ commutes with every letter and is a scalar by
    Lemma~\ref{lem:asym_normal_intertwiner}.
\end{proof}

\begin{corollary}[Two compressions with different flags in matrices]
    \label{cor:asym_uniq_different_flags_matrices}
    \lean{MPSTensor.MultiBlockCompression.leftRow_mul_eq_mul_leftRow_mul,
        MPSTensor.MultiBlockCompression.exists_rightCol_sub_mul_eq_fromCols_mul}
    \leanok
    \uses{thm:asym_uniq_different_flags, lem:asym_uniq_rank_test}
    In the setting of Theorem~\ref{thm:asym_uniq_different_flags}, with the matrices $G_{\leq k}$
    and $\widetilde G_{\leq k'}$ of Lemma~\ref{lem:asym_uniq_rank_test}: for every matrix $Y_U$ whose
    columns lie in $U$, $\widetilde W_{b'}Y_U=XW_bY_U$; and there is a matrix $R$ with
    $\widetilde V_{b'}-V_bX^{-1}=(G_{\leq k-1}\ \widetilde G_{\leq\pi(k)-1})\,R$.
\end{corollary}

\begin{proof}
    \leanok
    \uses{thm:asym_uniq_different_flags, lem:asym_uniq_rank_test}
    The first identity is the identity $\widetilde W_{b'}x=XW_bx$ of
    Theorem~\ref{thm:asym_uniq_different_flags} applied to every column of $Y_U$. By the same
    theorem every column of $\widetilde V_{b'}-V_bX^{-1}$ lies in
    $\mathcal N=M_{k-1}+\widetilde M_{\pi(k)-1}$, which is the column span of
    $(G_{\leq k-1}\ \widetilde G_{\leq\pi(k)-1})$ by Lemma~\ref{lem:asym_uniq_rank_test}.
\end{proof}

\begin{theorem}[A basis adapted to both flags]
    \label{thm:asym_uniq_common_basis}
    \lean{MPSTensor.MultiBlockCompression.matchedCompl,
        MPSTensor.MultiBlockCompression.finrank_matchedCompl,
        MPSTensor.MultiBlockCompression.flag_eq_iSup_matchedCompl,
        MPSTensor.MultiBlockCompression.flag_eq_iSup_matchedCompl',
        MPSTensor.MultiBlockCompression.iSupIndep_matchedCompl,
        MPSTensor.MultiBlockCompression.commonCoord,
        MPSTensor.MultiBlockCompression.blockTriangular_unipotentFactor,
        MPSTensor.MultiBlockCompression.blockDiag'_unipotentFactor,
        MPSTensor.MultiBlockCompression.matchingPermMatrix,
        MPSTensor.MultiBlockCompression.gaugeMatrix_mul_gaugeMatrixInv_eq}
    \leanok
    \uses{thm:asym_uniq_different_flags, prop:asym_uniq_matching}
    In the setting of Theorem~\ref{thm:asym_uniq_different_flags}, for two compressions of $B$ onto
    families $(C_s)_{s\in S}$ and $(\widetilde C_t)_{t\in\widetilde S}$ of normal tensors of
    positive bond dimension, let $U_k=M_k\cap\widetilde M_{\pi(k)}$,
    $K_k=M_{k-1}\cap\widetilde M_{\pi(k)}$, and let $L_k$ be a complement of $K_k$ in $U_k$. Then
    $\C^{D_B}=\bigoplus_kL_k$, $M_k=\bigoplus_{j\leq k}L_j$ and $\widetilde M_{k'}=\bigoplus_{\pi(j)\leq k'}L_j$,
    so the union of bases of the $L_j$ is adapted to both flags; and with $\Gamma$ the change of
    basis reading the coordinates of $L_j$ through $W_{\sigma(j)}$, one has $G=\mathsf A\Gamma$ and
    $\widetilde G=\widetilde{\mathsf A}P_X\Gamma$, where $\mathsf A$ and $\widetilde{\mathsf A}$ are block
    upper triangular with identity diagonal blocks in the orders $\sigma$ and $\widetilde\sigma$ and
    $P_X$ is the block permutation $j\mapsto\pi(j)$ composed with the intertwiners of
    Theorem~\ref{thm:asym_uniq_different_flags}; hence $\widetilde GG^{-1}=\widetilde{\mathsf A}P_X\mathsf A^{-1}$.
\end{theorem}

\begin{proof}
    \leanok
    \uses{thm:asym_uniq_different_flags, prop:asym_uniq_matching, lem:asym_uniq_maps}
    $U_k/K_k$ is isomorphic to $M_k/M_{k-1}$, so $\dim L_k$ is the size of the slot at $k$, and
    $M_k=M_{k-1}+U_k=M_{k-1}+K_k+L_k=M_{k-1}\oplus L_k$ since $K_k\subseteq M_{k-1}$ and $L_k\cap M_{k-1}\subseteq K_k$;
    induction on $k$, and the same for the second flag through $\pi$. $\Gamma$ is invertible because
    $W_{\sigma(j)}$ is injective on $L_j$. Both $G$ and $\Gamma$ map $M_k$ onto $E_k$, so $G\Gamma^{-1}$ is
    block upper triangular, with identity diagonal blocks because both read the step through
    $W_{\sigma(k)}$; reading the $L_j$ through $\widetilde W$ instead gives $X_j$ times the reading
    through $W$, by Theorem~\ref{thm:asym_uniq_different_flags}, which is $P_X\Gamma$.
\end{proof}

\begin{theorem}[Unipotent changes of two compressions]
    \label{thm:asym_uniq_unipotent_change}
    \lean{MPSTensor.MultiBlockCompression.flag_conjTriangular_and_unipotentFactor_conjTriangular,
        MPSTensor.MultiBlockCompression.matchedCompl_conjTriangular,
        MPSTensor.MultiBlockCompression.commonGauge_conjTriangular,
        MPSTensor.MultiBlockCompression.commonGauge'_conjTriangular,
        MPSTensor.MultiBlockCompression.matchingPermMatrix_conjTriangular,
        MPSTensor.MultiBlockCompression.exists_conjTriangular_unipotentFactor_eq}
    \leanok
    \uses{thm:asym_uniq_common_basis}
    In the setting of Theorem~\ref{thm:asym_uniq_common_basis}, let $u$ and $\widetilde u$ be block
    upper triangular with identity diagonal blocks in the orders $\sigma$ and $\widetilde\sigma$.
    Then $(\sigma,uG)$ and $(\widetilde\sigma,\widetilde u\widetilde G)$ are compressions with the
    same flags, the same complements $L_k$, the same change of basis $\Gamma$ and the same block
    permutation $P_X$, and with the unipotent factors $u\mathsf A$ and $\widetilde u\widetilde{\mathsf A}$.
    Conversely, for any matrices $\mathsf A_1$ and $\widetilde{\mathsf A}_1$ block upper triangular
    with identity diagonal blocks in the orders $\sigma$ and $\widetilde\sigma$ there are such $u$
    and $\widetilde u$ for which the factors are $\mathsf A_1$ and $\widetilde{\mathsf A}_1$: every
    pair of unipotent factors occurs.
\end{theorem}

\begin{proof}
    \leanok
    \uses{thm:asym_uniq_common_basis, thm:asym_uniq_same_flag}
    By Theorem~\ref{thm:asym_uniq_same_flag} the two changes keep the flags, hence the matching, the
    subspaces $U_k$ and $K_k$ and the complements $L_k$. The row of $uG$ at the slot of position $j$
    is $W_{\sigma(j)}+\sum_{c\text{ after }j}u_{jc}W_c$, and every $W_c$ with $c$ after $j$ vanishes
    on $L_j\subseteq M_j$, so $\Gamma$ is unchanged; likewise for $\widetilde u\widetilde G$, so $P_X$
    is unchanged. Then $uG=u\mathsf A\Gamma$ and $\widetilde u\widetilde G=\widetilde u\widetilde{\mathsf A}P_X\Gamma$.
    For the converse take $u=\mathsf A_1\mathsf A^{-1}$ and $\widetilde u=\widetilde{\mathsf A}_1\widetilde{\mathsf A}^{-1}$.
\end{proof}

\begin{theorem}[The set of all compressions]
    \label{thm:asym_uniq_all}
    \lean{MPSTensor.MultiBlockCompression.flag_withGauge_eq_iff_exists_eq_conjTriangularGauge,
        MPSTensor.FlagData.exists_multiBlockCompression_flag_eq,
        MPSTensor.FlagData.toMultiBlockCompression_ord_symm}
    \leanok
    \uses{prop:asym_uniq_dictionary, thm:asym_uniq_same_flag}
    Every labelled invariant flag of $B$ is the flag of a compression, and two compressions with the
    same order have the same flag exactly when they differ by a block upper triangular change of
    basis with nonzero scalar diagonal blocks. The set of all compressions of $B$ onto a family is
    therefore the disjoint union, over the labelled invariant flags of $B$, of one orbit of that group each.
\end{theorem}

\begin{proof}
    \leanok
    \uses{prop:asym_uniq_dictionary, thm:asym_uniq_same_flag}
    Proposition~\ref{prop:asym_uniq_dictionary} and Theorem~\ref{thm:asym_uniq_same_flag}.
\end{proof}

\begin{theorem}[The set of all reductions]
    \label{thm:asym_uniq_all_reductions}
    \lean{MPSTensor.MultiBlockCompression.exists_z_eq_iff_exists_flagData,
        MPSTensor.MultiBlockCompression.existsUnique_conjTriangular,
        MPSTensor.FlagData.toMultiBlockCompression_toFlagData,
        MPSTensor.FlagData.exists_rightInverse_toFlagData}
    \leanok
    \uses{prop:asym_uniq_dictionary, thm:asym_uniq_torsor}
    The tensor $B$ has a compression onto a family $(C_s)_{s\in S}$ with $z$ zero slots exactly when
    the word-algebra module of $B$ has a labelled invariant flag over $(S,C)$ with $z$ zero steps;
    this needs no hypothesis on the blocks. Let the blocks be normal of positive bond dimension,
    and choose for every labelled invariant flag $F$ a compression $(\sigma_F,G_F)$ whose labelled
    flag is $F$; such a choice exists. Then every compression of $B$ onto $(C_s)$ is
    $(\sigma_F,hG_F)$ for exactly one pair $(F,h)$ with $h\in T_{\sigma_F}$, and $F$ is its labelled
    flag. The compressions of $B$ onto $(C_s)$ are therefore the disjoint union, over the labelled
    invariant flags $F$, of the $T_{\sigma_F}$-orbits of $(\sigma_F,G_F)$.
\end{theorem}

\begin{proof}
    \leanok
    \uses{prop:asym_uniq_dictionary, thm:asym_uniq_torsor}
    The first statement is Proposition~\ref{prop:asym_uniq_dictionary} in both directions, and the
    compression constructed there from a labelled flag has that labelled flag, which gives the
    choice $F\mapsto(\sigma_F,G_F)$. Given a compression, $F$ must be its labelled flag; it then has
    the number of zero slots, the order and the flag of $(\sigma_F,G_F)$, and
    Theorem~\ref{thm:asym_uniq_torsor} gives exactly one $h$.
\end{proof}

\subsection{Long words and the one-block case}

\begin{lemma}[Long words in a triangular form]
    \label{lem:asym_uniq_long_words}
    \lean{MPSTensor.MultiBlockCompression.conjMatrix_evalWord_eq_zero_of_forall_notMem,
        MPSTensor.MultiBlockCompression.conjMatrix_evalWord_eq_zero_of_maxSlotPosition_lt,
        MPSTensor.MultiBlockCompression.conjMatrix_evalWord_eq_zero_of_lt_minSlotPosition,
        MPSTensor.MultiBlockCompression.evalWord_mulVec_mem_flag_of_forall_le,
        MPSTensor.MultiBlockCompression.evalWord_mulVec_eq_zero_of_mem_flag_of_forall_ge}
    \leanok
    \uses{thm:asym_multiblock}
    For a compression with $z$ zero slots and at least one block slot, and a word $w$ with
    $|w|\geq z$, the block $(GB^wG^{-1})_{pq}$ vanishes unless some block slot sits at a position
    $m$ with $p\leq m\leq q$. In particular $B^w$ maps $\C^{D_B}$ into $M_{k_{\max}}$ and kills
    $M_{k_{\min}-1}$, where $k_{\min}$ and $k_{\max}$ are the first and last positions of block slots.
\end{lemma}

\begin{proof}
    \leanok
    The block $(p,q)$ of a product of $n$ block upper triangular matrices is a sum over
    non-decreasing chains of positions from $p$ to $q$. If no block slot lies between $p$ and $q$,
    every chain runs through at most $z$ positions, so for $n\geq z$ it has an equal step whose
    factor is the diagonal block of a zero slot, which is zero.
\end{proof}

\begin{corollary}[Bare uniqueness at the ends of the flag; the one-block theorem]
    \label{cor:asym_uniq_one_block}
    \lean{MPSTensor.MultiBlockCompression.labelMatching_eq_of_unique,
        MPSTensor.MultiBlockCompression.exists_left_eq_smul_of_last,
        MPSTensor.MultiBlockCompression.exists_right_eq_inv_smul_of_first,
        MPSTensor.MultiBlockCompression.exists_left_mul_evalWord_eq_smul_of_last,
        MPSTensor.MultiBlockCompression.exists_evalWord_mul_right_eq_inv_smul_of_first,
        MPSTensor.MultiBlockCompression.exists_smul_single_block,
        MultiBlockExamples.SharpnessOneBlock.not_exists_left_eq_smul,
        MultiBlockExamples.SharpnessOneBlock.left_mul_evalWord_eq,
        MultiBlockExamples.SharpnessOneBlock.not_exists_left_mul_eq_smul_two,
        MultiBlockExamples.SharpnessOneBlock.left_mul_evalWord_eq_two,
        MultiBlockExamples.SharpnessOneBlock.exists_rowOnes_mul_pow_eq_smul_iff,
        MultiBlockExamples.SharpnessOneBlock.exists_length_eq_sub_one_not_exists_smul}
    \leanok
    \uses{thm:asym_uniq_different_flags, lem:asym_uniq_long_words, thm:asym_uniq_blocks}
    Let $C_s$ occur once in the family up to gauge. If $s$ is the last slot of both compressions
    then $\widetilde W_s=\lambda W_s$ as matrices, for a nonzero scalar $\lambda$; if first in both,
    $\widetilde V_s=\lambda^{-1}V_s$. If $s$ is the last block slot of both, then
    $\widetilde W_sB^w=\lambda W_sB^w$ for $|w|\geq z$; if the first block slot of both,
    $B^w\widetilde V_s=\lambda^{-1}B^wV_s$ for $|w|\geq z$. In particular, for a single block
    ($S=\{s\}$, $z=D_B-D_s$), there is $\lambda\neq0$ with $\widetilde W_sB^w=\lambda W_sB^w$ and
    $B^w\widetilde V_s=\lambda^{-1}B^wV_s$ for every word with $|w|\geq D_B-D_s$; for every $z\geq1$
    this bound cannot be lowered for the first identity.
\end{corollary}

\begin{proof}
    \leanok
    \uses{thm:asym_uniq_different_flags, lem:asym_uniq_long_words}
    Since $C_s$ occurs once, the matching pairs the two $s$-steps. If both are last, $U$ is the
    whole space; if both are first, $\mathcal N=0$. For the identities with words attached,
    Lemma~\ref{lem:asym_uniq_long_words} gives $B^w\C^{D_B}\subseteq U$ and $B^w\mathcal N=0$ for
    $|w|\geq z$, and Theorem~\ref{thm:asym_uniq_different_flags} applies to every $B^wx$. For the
    sharpness, the tensor with one letter acting on $\C^{1+z}$ by $e_0\mapsto e_0$, $e_j\mapsto e_{j-1}$
    has the compressions with the block first (identity gauge, $W=e_0^{\mathsf T}$) and with the
    block last (basis $e_{j-1}-e_j$, $e_0$; $\widetilde W=e_0^{\mathsf T}+\cdots+e_z^{\mathsf T}$), and
    $WB^n=e_0^{\mathsf T}+\cdots+e_{\min(n,z)}^{\mathsf T}$ is proportional to $\widetilde WB^n=\widetilde W$
    exactly when $n\geq z$.
\end{proof}

\subsection{Split compressions}

\begin{proposition}[Uniqueness of split compressions, bare form]
    \label{prop:asym_uniq_split}
    \lean{MPSTensor.MultiBlockCompression.remainder_eq_zero_iff_conjMatrix_eq_blockDiagonal',
        MPSTensor.MultiBlockCompression.blockEntry_intertwines,
        MPSTensor.MultiBlockCompression.blockEntry_eq_zero_of_not_isGaugeRelated,
        MPSTensor.MultiBlockCompression.blockEntry_inl_inr_eq_zero,
        MPSTensor.MultiBlockCompression.blockEntry_inr_inl_eq_zero,
        MPSTensor.MultiBlockCompression.left_eq_smul_of_split_of_pairwise_not_isGaugeRelated,
        MPSTensor.MultiBlockCompression.right_eq_smul_of_split_of_pairwise_not_isGaugeRelated,
        MPSTensor.MultiBlockCompression.splitScalar_ne_zero,
        MPSTensor.MultiBlockCompression.exists_left_eq_smul_and_right_eq_inv_smul_of_split,
        MPSTensor.MultiBlockCompression.classMatrix,
        MPSTensor.MultiBlockCompression.blockEntry_relativeGauge_eq_classMatrix_smul,
        MPSTensor.MultiBlockCompression.isUnit_classMatrix,
        MPSTensor.MultiBlockCompression.zeroSlotBlock_relativeGauge_mul_zeroSlotBlock_relativeGauge,
        MPSTensor.MultiBlockCompression.zeroSlotBlock_relativeGauge_mul_eq_one_and_mul_eq_one}
    \leanok
    \uses{thm:asym_multiblock, lem:asym_normal_intertwiner, thm:asym_zipper_uniqueness}
    Let two compressions of $B$ onto the same family of normal blocks both have zero remainder,
    $B^i=\sum_bV_bC_b^iW_b$. Then every block of $h=\widetilde GG^{-1}$ between two slots whose
    diagonal blocks are not gauge related, or between a block slot and a zero slot, vanishes;
    within one gauge class, with chosen intertwiners $Y_t$ from a representative $C_{s_0}$ to
    each member $C_t$, the blocks are $h_{t't}=\mu_{t't}\,Y_{t'}Y_t^{-1}$ for an invertible
    coefficient matrix $\mu$ indexed by the class; and the block of $h$ on the zero slots is
    invertible. If no block of the family is repeated up to gauge, then $\widetilde W_s=\lambda_sW_s$
    and $\widetilde V_s=\lambda_s^{-1}V_s$ as matrices for nonzero scalars $\lambda_s$, with no words
    attached. This is the bare form of Theorem~\ref{thm:asym_zipper_uniqueness} for the
    compressions of Theorem~\ref{thm:asym_multiblock}.
\end{proposition}

\begin{proof}
    \leanok
    \uses{lem:asym_normal_intertwiner}
    In the two coordinate systems the letters are block diagonal, so $h$ intertwines two block
    diagonal tensors and each block $h_{b'b}$ intertwines the diagonal block at $b$ with that at
    $b'$; Lemma~\ref{lem:asym_normal_intertwiner} makes it zero unless the two are gauge related,
    and between a block slot and a zero slot a nonzero intertwiner would give a nonzero vector
    killed by every letter of a normal tensor. Without repetition the row of $h$ at $s$ has the
    single nonzero block $\lambda_s1$, and likewise the column of $h^{-1}$, which intertwines in the
    reverse direction, with $\lambda_s\lambda'_s=1$ from $\widetilde W_s\widetilde V_s=1$.
\end{proof}
